\section{Transport of the ordinary source}
\label{tr:section}

The deterministic matrix estimate of \texttt{ward\_heat.tex} (subsection ``Deterministic matrices in the original invariant frame'') extends to a particular
field-dependent, gauge-covariant nonlocal matrix. The extension pays its first
and second transport derivatives under the original Wilson law. Its conditional
mean has an actual thin-separator gain. The construction reproduces the
degree-two force-square source at the identity, while leaving an explicit
nonlinear remainder. The remainder is not bounded extensively here.

Throughout this section the group is a fixed \(SU(N)\). Constants may depend on
\(N\), the original derivative convention and bounded original plaquette
incidence, but not on the torus side \(n\), the marks or \(\beta\ge1\). Use
\[
 \dd\mu=Z_\beta^{-1}e^{-\beta S}\dd\lambda,
 \qquad F_i=X_iS,\qquad H_{ij}=X_iF_j,
\]
where \(i=(b,A)\) is an original oriented link and color index. The derivative
\(X_b^A\) differentiates \(U_b\mapsto e^{tT_A}U_b\). In this convention the
force color frame is based at the origin \(o(b)\). Derivatives on different
links commute; derivatives on the same link retain their Lie brackets.
The matrix \(H\) is the ordered Hessian. Bounded determining incidence gives
\begin{equation}
 \norm H_{\rm op}\le C
 \label{tr:hessian}
\end{equation}
pointwise, including with arbitrary fixed unitary boundary values. No positive
lower Hessian bound is assumed.

The three original plane matchings are
\(\{12,34\}\), \(\{13,24\}\), and \(\{14,23\}\), with actions
\(S_1,S_2,S_3\). For an original link each matching contributes its complete
directional force. Thus the unchanged sources are
\[
 D[f]=\sum_{b,k}f_b\abs{\nabla_bS_k}^2,\qquad
 r_f=\beta D[f]-\ell[f]/3,\qquad
 h_f=\beta\bigl(D[f]-P[f]/3\bigr),\qquad M_f=\sum_bf_b^2.
\]
All marks are real and deterministic. All affected original roots, Casimir
factors and pointing conventions remain in these sums.

\subsection{An upper divergence estimate with the original Lie brackets}

For a smooth real vector field \(V=\sum_iV_iX_i\), define
\[
 \delta_\mu V=-\operatorname{div}_\mu V
   =\beta\sum_iF_iV_i-\sum_iX_iV_i.
\]
Compact integration by parts gives
\(\E((\delta_\mu V)G)=\E\sum_iV_iX_iG\).

\begin{lemma}[Upper divergence and own-link force estimates]
\label{tr:divergence}
For \(\beta\ge1\),
\begin{equation}
 \norm{\delta_\mu V}_2^2
 \le C\E\left(\sum_{i,j}\abs{X_jV_i}^2+\beta\abs V^2\right).
 \label{tr:divupper}
\end{equation}
If each coefficient \(W_j\) is independent of the entire original input link
of \(j\), then
\begin{equation}
 \E\abs{\sum_jW_jF_j}^2
 \le C\beta^{-1}\E\abs W^2
       +C\beta^{-2}\E\sum_{j,k}\abs{X_kW_j}^2.
 \label{tr:row}
\end{equation}
These are finite-volume upper estimates. They require neither a spectral gap
uniform in volume nor convexity of the Wilson potential.
\end{lemma}

\begin{proof}
Use \(X_i^*=\beta F_i-X_i\). A first integration by parts, followed by a
second one in the term containing \(F_j\), gives the exact identity
\begin{align}
 \E(\delta_\mu V)^2
 ={}&\beta\E\sum_{i,j}V_iV_jX_iF_j
    +\E\sum_{i,j}(X_jV_i)(X_iV_j)\notag\\
   &+\E\sum_{i,j}V_i[X_j,X_i]V_j.
 \label{tr:dividentity}
\end{align}
Indeed, differentiate \(\delta_\mu V\) against \(V_iX_i\), then integrate
\(\beta F_jV_iX_iV_j\) in \(j\). The two ordered second derivatives differ
by precisely the displayed bracket. For indices on different links this
bracket vanishes. For indices on the same link its fixed structure constants,
followed by Cauchy's inequality in that finite color block, bound its sum by
\(C\E(\abs V^2+\sum_{i,j}\abs{X_jV_i}^2)\).
Equation~\eqref{tr:hessian} bounds the first term. This proves
\eqref{tr:divupper}.

Under the own-link hypothesis, \(\sum_jX_jW_j=0\), so applying
\eqref{tr:divupper} to \(W\) already proves \eqref{tr:row}. An additional
direct calculation checks the coefficient responses:
\begin{equation}
 \beta^2\E(W\cdot F)^2
 =\beta\E\sum_{j,k}W_jW_kX_jF_k
       +\E\sum_{j,k}(X_kW_j)(X_jW_k).
 \label{tr:rowexact}
\end{equation}
Integrating first in \(j\) and then in \(k\) produces a possible extra term
\(\E\sum_{j,k}W_jX_kX_jW_k\). If the links differ, the derivatives commute
and \(X_kW_k=0\). If they are the same, \(X_jW_k=0\) by independence of the
entire link. Thus the extra term is identically zero. Equation
\eqref{tr:rowexact} and Cauchy's inequality give the stated bound again.
This is functional independence from a coordinate, not independence of random
links under \(\mu\).
\end{proof}

\subsection{A quantitative theorem for endpoint-independent moving matrices}

For a smooth real symmetric component matrix \(Q(U)\), put
\begin{align*}
 K_0(Q)&=\E\sum_{i,j}\abs{Q_{ij}}^2,\\
 K_1(Q)&=\E\sum_{k,i,j}\abs{X_kQ_{ij}}^2,\\
 K_2(Q)&=\E\sum_{l,k,i,j}\abs{X_lX_kQ_{ij}}^2.
\end{align*}
Pointwise upper envelopes for these sums may be used in place of the
expectations. Say that \(Q\) is endpoint-independent if \(Q_{ij}\) depends on
neither the entire link of \(i\) nor the entire link of \(j\).

\begin{theorem}[Moving matrix Ward estimate]
\label{tr:moving}
For an endpoint-independent \(Q\), set
\begin{equation}
 V_Q=QF,\qquad
 a_Q=\sum_{i,j}Q_{ij}X_iF_j,\qquad
 T_Q=\delta_\mu(QF)=\beta F^TQF-a_Q.
 \label{tr:TQ}
\end{equation}
Then \(\E T_Q=0\), and
\begin{align}
 \E\abs{QF}^2
 &\le C\bigl(K_0/\beta+K_1/\beta^2\bigr),
 \label{tr:QF}\\
 \Var_\mu T_Q
 &\le C\bigl(K_0+K_1/\beta+K_2/\beta^2\bigr),
 \label{tr:movingvariance}\\
 \inner{T_Q}{L_\mu^{-1}T_Q}_\mu
 &\le C\bigl(K_0/\beta+K_1/\beta^2\bigr).
 \label{tr:movinginverse}
\end{align}
Here \(L_\mu=-\sum_iX_i^2+\beta\sum_iF_iX_i\) is the canonical original
generator, and its inverse is on mean-zero functions. The same assertions
hold with arbitrary fixed unitary boundaries when the matrix hypotheses
and bounded original incidence hold.
\end{theorem}

\begin{proof}
Endpoint independence gives \(X_iQ_{ij}=0\) for every \(i,j\), proving the
exact divergence formula in \eqref{tr:TQ}. No moving coefficient has been
frozen. Apply \eqref{tr:row} to every input row \(j\mapsto Q_{ij}\) and sum
over \(i\); this proves \eqref{tr:QF}.

For each \(k,i\), the differentiated row \(j\mapsto X_kQ_{ij}\) is still
independent of the entire input link of \(j\). A derivative on another link
preserves that independence, while a derivative on the input link annihilates
the coefficient. Consequently a second use of the row estimate gives
\[
 \E\sum_{k,i}\abs{\sum_j(X_kQ_{ij})F_j}^2
 \le C\bigl(K_1/\beta+K_2/\beta^2\bigr).
\]
Since
\[
 X_k(QF)_i=\sum_j(X_kQ_{ij})F_j+\sum_jQ_{ij}H_{kj},
\]
the ordered Hessian bound yields
\begin{equation}
 \E\sum_{k,i}\abs{X_k(QF)_i}^2
 \le C\bigl(K_0+K_1/\beta+K_2/\beta^2\bigr).
 \label{tr:gradV}
\end{equation}
Use Lemma~\ref{tr:divergence} with \(V=QF\). This proves
\eqref{tr:movingvariance}, including the contributions from all first and
second coefficient derivatives and the same-link brackets.

Finally \(\E(T_QG)=\E(QF)\cdot\nabla G\). Cauchy's inequality and the
variational inverse-form identity bound \(\inner{T_Q}{L_\mu^{-1}T_Q}\)
by \(\E\abs{QF}^2\). This proves \eqref{tr:movinginverse}; no uniform
spectral gap has been introduced in using the inverse.
\end{proof}

\subsection{A flat reference with a uniform weighted Hilbert--Schmidt bound}

This subsection constructs deterministic matrices on the periodic \(n^4\)
torus, \(n\ge3\). Let \(H_k^0\) be the identity Hessian of \(S_k\) and
\(H_0^0=\sum_kH_k^0\). With
\[
 q_\mu(p)=e^{ip_\mu}-1,\qquad s(p)=\sum_\mu\abs{q_\mu(p)}^2,
\]
the original common quadratic normalization \(\kappa>0\) gives
\[
 H_0^0(p)=\kappa\bigl(s(p)I-q(p)q(p)^*\bigr).
\]
Each \(H_k^0\) is nonnegative, annihilates \(q\), and is bounded above by
\(H_0^0\). Define
\begin{equation}
 A_k(p)=\frac{H_k^0(p)}{\kappa s(p)}\quad(p\ne0),\qquad A_k(0)=0.
 \label{tr:Ak}
\end{equation}
Equivalently \(A_k=H_k^0(H_0^0)^+\); hence
\(\norm{A_k}_{\rm op}\le1\) and \(H_k^0=A_kH_0^0\).
These are real scalar matrices in link space, tensored with the identity in
color. The zero-mode convention is made only in this deterministic reference;
it does not remove any modes or holonomies from the compact measure.

For real marks let
\begin{equation}
 Q_f^0=\sum_{k=1}^3A_k^*D_fA_k,\qquad (D_f)_{bc}=f_b\delta_{bc},
 \qquad q^f_{bc}=(Q_f^0)_{bc}.
 \label{tr:Qflat}
\end{equation}
The matrix is real symmetric. The Hilbert--Schmidt ideal inequality directly
gives \(\norm{Q_f^0}_{\HS}^2\le CM_f\), including signed marks.

\begin{lemma}[Periodic weighted matrix estimate]
\label{tr:weightedHS}
With \(d_T\) the torus graph distance between link origins,
\begin{equation}
 \sum_{b,c}\bigl(1+d_T(o(b),o(c))^2\bigr)\abs{q^f_{bc}}^2
 \le CM_f.
 \label{tr:weighted}
\end{equation}
The constant is uniform in \(n\); there is no absolute-kernel logarithmic
factor.
\end{lemma}

\begin{proof}
For each coordinate \(\nu\), use the periodic multiplication
\[
 \phi_\nu(x)=\frac{n}{2\pi}e^{2\pi ix_\nu/n}.
\]
The sum of \(\abs{\phi_\nu(x)-\phi_\nu(y)}^2\) over \(\nu\) is comparable
to \(d_T(x,y)^2\). Uniform row and column estimates for each \(A_k\) are
\begin{equation}
 \sup_x\sum_y\abs{\phi_\nu(x)-\phi_\nu(y)}^2
                   \norm{A_k(x,y)}_{\HS}^2\le C.
 \label{tr:weightedrow}
\end{equation}
Here the finite matrix in \(x,y\) includes the four link orientations.
To verify this, set \(h=2\pi/n\). Parseval converts the row sum, up to a
fixed factor, to
\begin{equation}
 \frac{n^2}{n^4}\sum_p
       \norm{A_k(p+he_\nu)-A_k(p)}_{\HS}^2.
 \label{tr:diffparseval}
\end{equation}
Away from a distance \(2h\) of zero and the shifted singular point, the
explicit order-zero symbol in \eqref{tr:Ak} satisfies
\(\norm{\nabla_pA_k}\le C/\operatorname{dist}(p,0)\).
Its finite difference is at most \(Ch/\operatorname{dist}(p,0)\).
In four dimensions the normalized sum of
\(\operatorname{dist}(p,0)^{-2}\) away from the origin is bounded: shell
counting reduces it to the integrable radial power \(r^{3-2}\).
The bounded number of modes within distance \(2h\) contributes at most
\(Cn^2/n^4=Cn^{-2}\) to \eqref{tr:diffparseval}, using
\(\norm{A_k}_{\rm op}\le1\). This proves \eqref{tr:weightedrow} and its
column counterpart.

The exact commutator is
\[
 [\phi,Q_f^0]=\sum_k\bigl([\phi,A_k^*]D_fA_k
                          +A_k^*D_f[\phi,A_k]\bigr).
\]
Multiplication by \(D_f\), \eqref{tr:weightedrow}, and the
Hilbert--Schmidt ideal inequality bound its squared norm by \(CM_f\).
For the adjoint term one uses
\([\phi,A_k^*]=-[\overline\phi,A_k]^*\); the complex multiplication
\(\phi\) is not being treated as self-adjoint. Summing in \(\nu\), then
adding the unweighted estimate, proves \eqref{tr:weighted}.
\end{proof}

\subsection{Endpoint-avoiding parallel transport}

For every unordered pair of original links \(b,c\), choose a deterministic
path \(\gamma_{bc}\) from \(o(b)\) to \(o(c)\) that avoids both underlying
links \(b,c\), and put \(\gamma_{cb}=\gamma_{bc}^{-1}\). If the origins agree,
choose the empty path. A shortest path encounters each forbidden edge at most
once. Replace each such occurrence by a three-edge elementary plaquette
detour avoiding both edges. In dimension four a perpendicular direction can
be chosen to avoid the other forbidden edge. At most two detours are needed,
so these paths have length
\begin{equation}
 L_{bc}\le d_T(o(b),o(c))+4
 \label{tr:pathlength}
\end{equation}
and uniformly bounded edge multiplicity. Keeping the detours as walks
preserves their chosen shortest-path homotopy. When \(b,c\) share an original
plaquette, choose its local arc and elementary detours in a fixed neighborhood.
The complete path family is chosen independently of the marks and of any
later random separator offset.

Define the actual compact-field matrix
\begin{equation}
 (Q_f(U))_{bA,cB}
    =q^f_{bc}\,[\Ad U(\gamma_{bc})]_{AB}.
 \label{tr:Qtransport}
\end{equation}
It is real symmetric and covariant in the two origin color frames. It is
independent of both endpoint links. Its transport paths can pass through
other links that will later be integrated; no such dependence is suppressed.

\begin{theorem}[Extensive actual-law transported Ward component]
\label{tr:ward}
For the periodic original Wilson law, \(n\ge3\) and \(\beta\ge1\), let
\(T_f=T_{Q_f}\) be defined by \eqref{tr:TQ} and
\eqref{tr:Qtransport}. Then
\begin{equation}
 \E T_f=0,\qquad
 \Var_\mu T_f\le CM_f,\qquad
 \inner{T_f}{L_\mu^{-1}T_f}_\mu\le C\beta^{-1}M_f,
 \qquad \E\abs{Q_fF}^2\le C\beta^{-1}M_f.
 \label{tr:wardbounds}
\end{equation}
These conclusions concern the constructed Ward component under the actual
compact law. They do not assert a bound for \(r_f\).
\end{theorem}

\begin{proof}
Each adjoint transporter is orthogonal. A derivative on one of the at most
\(L_{bc}\) path links is a sum of a uniformly bounded number of insertions
of fixed adjoint generators. The edge multiplicity is bounded. Thus,
pointwise on the complete compact link product,
\begin{align*}
 \sum_k\norm{X_k\Ad U(\gamma_{bc})}_{\HS}^2&\le CL_{bc},\\
 \sum_{l,k}\norm{X_lX_k\Ad U(\gamma_{bc})}_{\HS}^2&\le CL_{bc}^2.
\end{align*}
All ordered second derivatives, including two derivatives on the same link,
are included. Multiplying by \(\abs{q^f_{bc}}^2\), summing over endpoints,
and using \eqref{tr:pathlength} and Lemma~\ref{tr:weightedHS}, gives
pointwise envelopes with
\begin{equation}
 K_0(Q_f)+K_1(Q_f)+K_2(Q_f)\le CM_f.
 \label{tr:Kbounds}
\end{equation}
Theorem~\ref{tr:moving} proves \eqref{tr:wardbounds}.
No pointwise operator norm of the transported matrix was estimated. In
particular, replacing signed Riesz cancellations by an absolute kernel sum,
which could cost \(\log n\), was unnecessary.
\end{proof}

\begin{proposition}[Complete ordinary contacts]
\label{tr:ordinarycontact}
For \(R_g=r_g\) or \(h_g\), put \(G_g=R_g/\beta\). Then
\begin{equation}
 \Cov_\mu(T_f,R_g)
   =\E\sum_{i,j}(Q_f)_{ij}X_jX_iG_g,
 \qquad
 \abs{\Cov_\mu(T_f,R_g)}\le C\sqrt{M_fM_g}.
 \label{tr:sourcecontact}
\end{equation}
The contact identity applies to every endpoint-independent moving matrix.
The bound uses its \(K_0\) estimate and requires no directional moment input.
\end{proposition}

\begin{proof}
The complete local source stencil and compactness give, pointwise,
\[
 \norm{\Hess_{\rm ord}G_g}_{\HS}^2\le CM_g.
\]
For \(r_g\), this includes the derivative of
\(G_g=D[g]-\ell[g]/(3\beta)\); \(\beta\ge1\) suffices. To check the bound,
each ordered Hessian entry reads boundedly many original source roots and
each source root contributes to boundedly many entries. Cauchy's inequality
on these finite incidence sets gives \(CM_g\).

Since \(\E T_f=0\), two original integrations by parts give
\[
 \E(T_fR_g)
 =\beta\E\sum_{i,j}(Q_f)_{ij}F_jX_iG_g
 =\E\sum_{i,j}(Q_f)_{ij}X_jX_iG_g.
\]
The coefficient derivative \(X_j(Q_f)_{ij}\) is exactly zero by input-link
independence. No ordered derivatives are commuted. Hilbert--Schmidt Cauchy
and \eqref{tr:Kbounds} prove the result. This argument also works with
arbitrary fixed unitary boundaries for matrices satisfying the same bounds.
\end{proof}
\subsection{Exact compression retains the raw marginal response}

Split the original links into integrated coordinates \(I\) and retained
coordinates \(B\). Let \(K\) be the complete compact conditional expectation,
with exact marginal
\begin{equation}
 \dd\nu_B=e^{-V_{\rm eff}}\dd\lambda_B,\qquad
 e^{-V_{\rm eff}(b)}
   =Z_\beta^{-1}\int e^{-\beta S(u_I,b)}\dd\lambda_I(u_I).
 \label{tr:rawmarginal}
\end{equation}
For a cube observation, this density is exactly
\(Z_\beta^{-1}e^{-\beta S_{\rm out}(b)}\prod_jZ_j(b)\).
For the original one-link layer, it retains its original \(\prod_sZ_s(b)\).
No induced action is replaced by a bare Wilson action.

Compact integration by parts gives
\begin{equation}
 KT_f=-\operatorname{div}_{\nu_B}\overline V_f,\qquad
 \overline V_f=K((Q_fF)_B).
 \label{tr:compression}
\end{equation}
Indeed, the integrated output divergences vanish and the retained output
divergence is the divergence in the true marginal. For each retained output
\(a\), every integrated input column satisfies
\[
 K((Q_f)_{ac}F_c)=\beta^{-1}K(X_c(Q_f)_{ac})=0,
 \qquad c\in I.
\]
The remaining matrix still depends on the integrated coordinates along its
paths. Put \(\overline Q_{ab}=K(Q_f)_{ab}\) for retained \(a,b\). The exact
retained vector field is
\begin{equation}
 (\overline V_f)_a
  =\beta^{-1}\left(
     \sum_{b\in B}\overline Q_{ab}X_bV_{\rm eff}
       -\sum_{b\in B}X_b\overline Q_{ab}\right).
 \label{tr:movingcompression}
\end{equation}
The minus sign and the second term are essential. To check them, differentiate
the original normalized conditional integral:
\[
 X_b\overline Q_{ab}
 =K(X_b(Q_f)_{ab})-\beta\Cov_K((Q_f)_{ab},F_b)
 =-\beta\Cov_K((Q_f)_{ab},F_b).
\]
Also \(KF_b=\beta^{-1}X_bV_{\rm eff}\). Solving the last display for
\(K((Q_f)_{ab}F_b)\) gives \eqref{tr:movingcompression}. Thus the
raw-normalizer covariance response remains in the compressed field, even
though the original endpoint coefficient derivatives vanish.

Conditional Jensen gives \(\Var_{\nu_B}(KT_f)\le CM_f\).
For the canonical generator of this same marginal, with its original product
metric,
\begin{equation}
 \inner{KT_f}{L_{\nu_B}^{-1}KT_f}_{\nu_B}
  \le\E_{\nu_B}\abs{\overline V_f}^2
  \le C\beta^{-1}M_f.
 \label{tr:compressedinverse}
\end{equation}
This is a generator-specific conclusion. Conditional Jensen also applies to
other complete positive observations, but it does not identify their history
generators with \(L_{\nu_B}\). Equations
\eqref{tr:compression}--\eqref{tr:movingcompression} integrate complete
compact coordinates. They do not differentiate a sharp selected boundary.
For a good/bad partition, the raw branches must first recombine to that
complete coordinate integral before these formulas are used.

\subsection{An actual thin-separator gain}

Let \(\tau\) be an independent packing offset, and let \(B_\tau\) be the
retained original link set. Suppose
\begin{equation}
 \sup_b\Pr_\tau(b\in B_\tau)\le p.
 \label{tr:thincondition}
\end{equation}
The paths defining \(Q_f\) are fixed independently of \(\tau\).
Let \(K_\tau\) integrate the complete complement of \(B_\tau\) with its
actual raw marginal. No product conditional law is needed in the following
theorem. Write \(I_B\) for the component projection onto retained links and
define
\[
 Q_f^B=I_BQ_fI_B,\qquad W_f^B=\delta_\mu(Q_f^BF).
\]

\begin{theorem}[Thin-separator mean gain in the actual compact law]
\label{tr:separator}
For every fixed offset,
\begin{equation}
 K_\tau T_f=K_\tau W_f^{B_\tau}.
 \label{tr:maskidentity}
\end{equation}
Moreover,
\begin{align}
 \Var_{\tau,\mu}(K_\tau T_f)&\le CpM_f,
 \label{tr:meanthin}\\
 \E_\tau\inner{K_\tau T_f}
       {L_{\nu_{B_\tau}}^{-1}K_\tau T_f}_{\nu_{B_\tau}}
   &\le Cp\beta^{-1}M_f.
 \label{tr:inversethin}
\end{align}
For the separated-cube packing with \(p\le C/R\), the first bound is an
actual \(R^{-1}\) gain. It is not a Gaussian finite-rank comparison.
\end{theorem}

\begin{proof}
In \eqref{tr:compression}, output divergences at deleted links integrate
to zero. For a retained output, every deleted input column has zero conditional
expectation by endpoint independence and compact integration by parts. The
surviving vector field is therefore exactly the one obtained from
\(Q_f^{B_\tau}\). This proves \eqref{tr:maskidentity}, while retaining
the dependence of the transport paths on all integrated links.

Once \(\tau\) is fixed, the endpoint mask has no field derivative. Every
squared coefficient and every first or second derivative contribution in
\(K_j(Q_f^{B_\tau})\) has a factor
\(\mathbf 1_{\{b,c\in B_\tau\}}\). This factor is bounded by
\(\mathbf 1_{\{b\in B_\tau\}}\). Average it against \(\tau\) before
summing the deterministic path envelopes used in
\eqref{tr:Kbounds}. The result is
\begin{equation}
 \E_\tau\bigl(K_0(Q_f^{B_\tau})+K_1(Q_f^{B_\tau})
                         +K_2(Q_f^{B_\tau})\bigr)\le CpM_f.
 \label{tr:maskedK}
\end{equation}
No interchange of a supremum in \(U\) with expectation is required: the
pointwise sums have the deterministic envelopes already proved.

Theorem~\ref{tr:moving}, conditional Jensen, and
\eqref{tr:maskidentity} give
\[
 \E_\tau\norm{K_\tau T_f}_2^2
 \le\E_\tau\norm{W_f^{B_\tau}}_2^2\le CpM_f.
\]
For every offset, \(\E_\mu T_f=0\). Thus there is no discarded
between-offset mean variance in \eqref{tr:meanthin}. Apply the variational
argument of \eqref{tr:compressedinverse} to the masked retained field and
use \eqref{tr:QF}, averaged as in \eqref{tr:maskedK}. This proves
\eqref{tr:inversethin} for the actual canonical marginal generators.
\end{proof}

The complete ordinary contacts have a sharper averaged masked estimate:
\begin{equation}
 \E_\tau\abs{\Cov_\mu(W_f^{B_\tau},R_g)}
       \le Cp\sqrt{M_fM_g},\qquad R_g=r_g\ \hbox{or}\ h_g.
 \label{tr:maskedcontact}
\end{equation}
To verify it, use the exact contact in
Proposition~\ref{tr:ordinarycontact}. Before integrating \(U\), bound its
absolute value by the sum of
\(\mathbf 1_{\{b,c\in B_\tau\}}\abs{q^f_{bc}}\) times the Frobenius
norm of the corresponding color block of
\(\Hess_{\rm ord}(R_g/\beta)\), with a fixed group factor from
\(\Ad U(\gamma_{bc})\). The Hessian and the complete marks are independent
of \(\tau\). Average the indicator first, obtaining \(p\), and then use
Cauchy's inequality in the full endpoint sum. This proves
\eqref{tr:maskedcontact} without a moment assumption or a conditional
centering statement.

\begin{corollary}[Projected ordinary covariance comparison]
\label{tr:cross}
Suppose the chosen complete coordinate observation also satisfies the
separate averaged internal estimate
\begin{equation}
 \E_\tau\E_\mu\Var(r_g\mid B_\tau)\le C M_g.
 \label{tr:internalinput}
\end{equation}
For the explicitly constructed separated-cube observation this is the
all-branch internal theorem proved in the accompanying construction, under
its stated geometric and original-moment hypotheses. It is not an
identification with R-II. Put
\[
 m_f=K_\tau r_f-\E_\mu r_f,\qquad
 w_f=K_\tau T_f,\qquad
 e_f=K_\tau(r_f-T_f)-\E_\mu r_f=m_f-w_f.
\]
Then, under the joint law of \(\tau\) and the actual retained coordinates,
\begin{align}
 \abs{\inner{m_g}{w_f}}
   &\le C(p+\sqrt p)\sqrt{M_fM_g},
 \label{tr:meancross}\\
 \abs{\Cov(m_f,m_g)-\Cov(e_f,e_g)}
   &\le C(p+\sqrt p)\sqrt{M_fM_g}.
 \label{tr:meancomparison}
\end{align}
In particular the change of squared variance for one mark is
\(O(R^{-1/2}M_f)\) when \(p\le C/R\).
\end{corollary}

\begin{proof}
Keep the actual total covariance identity, using
\eqref{tr:maskidentity}:
\begin{align*}
 \inner{m_g}{w_f}
 ={}&\E_\tau\Cov_\mu(r_g,W_f^{B_\tau})\\
    &-\E_\tau\E_\mu\Cov(r_g,W_f^{B_\tau}\mid B_\tau).
\end{align*}
The first term is bounded by \eqref{tr:maskedcontact}. Conditional Cauchy,
then Cauchy in the actual outside law and \(\tau\), bounds the second by
\[
 \left(\E_\tau\E\Var(r_g\mid B_\tau)\right)^{1/2}
 \left(\E_\tau\E\Var(W_f^{B_\tau}\mid B_\tau)\right)^{1/2}
 \le C\sqrt p\sqrt{M_fM_g}.
\]
Here the first factor uses \eqref{tr:internalinput}; the second uses the
full original-law masked Ward variance, not a fixed-bath estimate.
This proves \eqref{tr:meancross}. Since \(e_f=m_f-w_f\), retain both
cross terms in the exact expansion
\[
 \Cov(m_f,m_g)-\Cov(e_f,e_g)
 =\Cov(m_f,w_g)+\Cov(w_f,m_g)-\Cov(w_f,w_g).
\]
The first two terms use \eqref{tr:meancross}, and the last uses
\eqref{tr:meanthin}. This proves \eqref{tr:meancomparison}.
\end{proof}

\subsection{The exact nonlinear residual and a paid contact estimate}

Define the following residual without changing the ordinary source:
\begin{equation}
 E_f=r_f-T_f
   =\beta\bigl[D[f]-F^TQ_f(U)F\bigr]+C_f,
 \qquad C_f=a_{Q_f}-\ell[f]/3.
 \label{tr:residual}
\end{equation}
Theorem~\ref{tr:ward} and Proposition~\ref{tr:ordinarycontact} give
\begin{equation}
 \abs{\Var_\mu r_f-\Var_\mu E_f}\le CM_f.
 \label{tr:fullcomparison}
\end{equation}
At the identity, the degree-two force-square jet cancels exactly:
\begin{equation}
 \sum_{b,k}f_b\abs{H_k^0A}_b^2
       =(H_0^0A)^TQ_f^0(H_0^0A).
 \label{tr:jet}
\end{equation}
This follows from \(H_k^0=A_kH_0^0\), with no replacement of the global
source by independent local kinetic pieces. Equation~\eqref{tr:jet} is
an identity-chart Taylor statement, not a global reconstruction theorem.

\begin{proposition}[Local contact remainder under the actual law]
\label{tr:contactremainder}
Use the local contractible path convention for endpoint pairs that share an
original plaquette. Let \(n\ge n_0\), where \(n_0\) is a fixed size large
enough for those bounded local patches to embed. In the original periodic
moment regime \eqref{rt:momentregime}, for every \(q\ge1\),
\begin{equation}
 \norm{C_f-C_f(I)}_{L^q(\mu)}
     \le Cq\beta^{-1}\norm f_1.
 \label{tr:contactLp}
\end{equation}
Consequently, for any complete observation \(K\),
\begin{equation}
 \Var_\mu C_f\le C\beta^{-2}\norm f_1^2,
 \qquad
 \Var(KC_f)\le C\beta^{-2}\norm f_1^2
       \le C\beta^{-2}\abs{\supp f}M_f.
 \label{tr:contactvariance}
\end{equation}
This is extensive for marks with \(\abs{\supp f}\le C\beta^2\). It is not
a volume-uniform extensive estimate for arbitrary marks.
\end{proposition}

\begin{proof}
The ordered Hessian block \(H_{bc}\) vanishes unless \(b,c\) belong to a
common original plaquette. For such a pair set
\[
 g_{bc}(U)=\sum_{A,B}[\Ad U(\gamma_{bc})]_{AB}
                         H_{(b,A),(c,B)}(U).
\]
This scalar is gauge invariant: the transporter and Hessian block transform
by the same orthogonal adjoint matrices at their two origins. By the local
path convention it is supported in a fixed contractible neighborhood.
At the identity every first link derivative of \(g_{bc}\) vanishes. Indeed,
global conjugation fixes the identity field, so that derivative would define
an invariant vector in the adjoint representation of \(\mathfrak{su}(N)\),
which has none.

Choose a fixed embedded local box containing the plaquette and its detours,
and a local spanning tree. Its remaining loop variables have plaquette
fillings of bounded area. Nonabelian plaquette reconstruction, with all
conjugating prefixes retained, gives a gauge in which their total squared
distance to identity is at most \(C\sum_{p\in\mathcal P_{bc}}s_p\).
Here \(\mathcal P_{bc}\) is a fixed-size set of original plaquettes. If this
sum is small, all the remaining local links lie in a common chart; Taylor's
theorem and the vanishing first derivative give
\begin{equation}
 \abs{g_{bc}(U)-g_{bc}(I)}
      \le C\sum_{p\in\mathcal P_{bc}}s_p.
 \label{tr:localcontact}
\end{equation}
On the compact complement, increase \(C\) using the positive lower threshold
for that sum. This proves \eqref{tr:localcontact} globally. This local
gauge is only an estimate of a gauge-invariant scalar; it does not change
the actual outside law or remove a global holonomy. The explicit
\(n\ge n_0\) hypothesis makes this fixed-patch argument literal.

Write \(b\sim c\) for pairs sharing a plaquette, including the diagonal.
This relation has bounded degree. From \eqref{tr:Qflat},
\begin{align*}
 \sum_{b\sim c}\abs{q^f_{bc}}
 &\le\sum_a\abs{f_a}\sum_k\sum_{b\sim c}
                           \abs{(A_k)_{ab}(A_k)_{ac}}\\
 &\le C\sum_a\abs{f_a}\sum_k\sum_b\abs{(A_k)_{ab}}^2
 \le C\norm f_1.
\end{align*}
The last step uses \(\norm{A_k}_{\rm op}\le1\).
Also \(\ell[f]-\ell[f](I)\) is exactly a bounded-incidence sum of original
plaquette deficits with the original Casimir coefficient: each occurrence
of \(\Delta_bs_p\) is \(C_F\Re\Tr U_p\).
The original local exponential moment gives
\(\norm{s_p}_{L^q(\mu)}\le Cq/\beta\).
Minkowski's inequality, \eqref{tr:localcontact}, and the last coefficient
sum prove \eqref{tr:contactLp}. Centering and conditional Jensen prove
\eqref{tr:contactvariance}. No spatial covariance cancellation follows
from this local moment argument; the factor \(\norm f_1^2\) is retained.
\end{proof}

\subsection{A flat holonomy calculation limits the Taylor cancellation}

\begin{proposition}[A nonzero quadratic defect at a flat toron]
\label{tr:toron}
For the transported matrix above, the force-square residual in
\eqref{tr:residual} need not have a vanishing quadratic jet at a flat
connection with nontrivial holonomy. This occurs already in \(SU(2)\) with
constant marks.
\end{proposition}

\begin{proof}
Take odd \(n\ge\max\{3,n_0\}\), \(f_b=1\), and the constant flat background
\[
 U_{x,1}=g=\begin{pmatrix}e^{i\theta/(2n)}&0\\0&e^{-i\theta/(2n)}\end{pmatrix},
 \qquad U_{x,\nu}=I\quad(\nu\ne1),
\]
where \(0<\abs\theta<\pi\). Every plaquette is identity, whereas the
first-direction holonomy \(g^n\) is nontrivial. On the charged two-dimensional
color plane, viewed as one complex coordinate, \(\Ad g\) acts by
\(e^{i\theta/n}\).

Choose the signed shortest displacement representatives
\(d_\mu\in\{-(n-1)/2,\ldots,(n-1)/2\}\), with elementary endpoint
detours homotopic to those paths. Because the connection is flat, the
transporter between endpoints separated by \(d\) acts on this color plane
by \(e^{i\theta d_1/n}\). The detours preserve this value.

For momenta \(p\) with \(p_2=p_3=p_4=0\), the flat matrices satisfy
\[
 Q_1^0(p)=\operatorname{diag}(0,1,1,1)\quad(p_1\ne0),
 \qquad Q_1^0(0)=0.
\]
Indeed, for each transverse direction \(\nu\), only the matching containing
the plane \((1,\nu)\) acts, and its \(A_k\) is the unit projection onto that
component. The translation-invariant kernel therefore obeys
\[
 \sum_{d_2,d_3,d_4}q^1_{(0,\nu),(d,\nu)}
     =\mathbf 1_{\{d_1=0\}}-1/n\qquad(\nu\ne1),
\]
while the off-diagonal output-direction sums vanish. On a spatially constant
charged perturbation \(A\) in link direction \(\nu=2\), the transported
matrix consequently acts by the scalar
\begin{equation}
 q_\theta=1-D_n(\theta),\qquad
 D_n(\theta)=\frac1n\sum_{d=-(n-1)/2}^{(n-1)/2}e^{i\theta d/n}
       =\frac{\sin(\theta/2)}{n\sin(\theta/(2n))}.
 \label{tr:toronmultiplier}
\end{equation}
The symmetric sum is real. Reversing the Fourier or transport sign leaves
this scalar unchanged.

The actual Hessian at this flat background is the covariant Maxwell Hessian.
On the same perturbation it satisfies
\[
 H_0^\theta A=H_{\{12,34\}}^\theta A
    =\kappa s_\theta A,\qquad
 s_\theta=\abs{e^{i\theta/n}-1}^2,
\]
and the other two matching Hessians annihilate \(A\). The total and each
directional force vanish at the background itself. Hence, along the original
compact perturbation \(U(\varepsilon)\) with tangent \(A\),
\begin{equation}
 D[1](U(\varepsilon))-F^TQ_1(U(\varepsilon))F
   =\varepsilon^2D_n(\theta)\norm{H_0^\theta A}_2^2
       +O(\varepsilon^3).
 \label{tr:torondefect}
\end{equation}
For small nonzero \(\theta\) the quadratic coefficient is nonzero and
\(D_n(\theta)\) tends to one as \(\theta\) tends to zero. Thus vanishing
background plaquette actions do not make this residual cubic around every
flat compact connection.
\end{proof}

This example does not refute an extensive covariance bound. The mode in
\eqref{tr:torondefect} is soft:
\(s_\theta=O(\theta^2/n^2)\). In its stated quadratic Gaussian reference,
the variance of the displayed energy scales as \(O(\theta^4/n^4)\).
That calculation only checks its scale; it is not a comparison with the
actual compact law. A nonlinear proof must accommodate these holonomy-dependent
soft modes or change the reconstruction while paying its new responses.

\subsection{What the quantitative estimates leave open}

For the constructed separated-cube observation, the internal ordinary
variance is bounded under its stated hypotheses, and
Corollary~\ref{tr:cross} compares the retained ordinary mean to the retained
mean of the exact residual \eqref{tr:residual}. A sufficient missing bound
is now explicitly
\begin{equation}
 \Var_{\tau,\mu}\bigl(K_\tau E_f\bigr)\le CM_f,
 \label{tr:missing}
\end{equation}
or a proved absorption estimate
\begin{equation}
 \Var_{\tau,\mu}\bigl(K_\tau E_f\bigr)
   \le\eta\Var_{\tau,\mu}\bigl(K_\tau r_f\bigr)+CM_f,
 \qquad \eta<1.
 \label{tr:missingabsorption}
\end{equation}
Neither estimate is established. The local contact estimate retains the
nonextensive \(\norm f_1^2\) factor. The force residual retains its nonlinear
curvature, transport and holonomy responses. Score integration by parts
bounds the Ward component and its contacts; repeating the variance comparison
without another analytic input leaves coefficient one and supplies no
\(\eta<1\).

These results hold at the original source-zero law and for the stated complete
coordinate observations. They prove neither bounds under arbitrary joint
source tilts nor mixed complex admission. Sharp-selector first \(L^1\)
surface estimates remain distinct from higher selector words and complex
domains. The original flat-block restriction
\(\beta\ge CR^{14}\log(2+R)\), all actual raw \(Z_s\) and \(Z_j\),
complete source partners and heat acceleration, current and anchor responses,
inherited history generators, reserve ownership and S94 remain in force.
Physical multiplication by \(c_\beta^{-1}a^{-2}\) multiplies every stated
variance bound by \(c_\beta^{-2}a^{-4}\). No physical-scale induction by
tiling and no complete Yang--Mills existence proof follow from this section.
